\chapter*{پیش‌گفتار}
به سبب رشد نرم‌افزار نوپای زی‌پرشین در ایران و تنوع و پیچیدگی کار نیاز به یک راهنمای کوتاه و جامع و به روز را در دانشکده خود احساس کردیم، چرا که تا آن‌جا که یافتیم به روزترین راهنما ترجمه‌ی دکتر امیدعلی به نام مروری نه چندان کوتاه بر تک بود که هر چند مجموعه‌ای جامع است اما به دلیل عدم بیان روش نصب و حجم زیاد مطالب به نظرمان ناکافی آمد، ضمن این‌که چندی‌ست که دیگر میک‌تک به کناری رفته و تک‌لایو جایگزین آن شده. از این رو سعی مان بر اینست که موجز و مفید از نصب تا تکمیل کار را به صورت عملی در این شبه رساله بیان کنیم. ضمن اینکه در همین راستا به معرفی قالب طراحی شده برای دانشکده ریاضی دانشگاه فردوسی را که نسخه اصلاح‌شده‌ای از روی نمونه‌های تهیه شده به وسیله آقایان خلیقی و دامن‌افشان را ارائه خواهیم داد.\\
باشد که مورد استفاده دوستانمان و رضای خدا واقع شود.